\section*{2 Preliminaries}
\paragraph{Notations}
For $r\in\mathbb N$, $[r]:=\{1,\ldots,r\}$. We use lowercase Greek alphabets like $\alpha,\beta$ for field constants, bold-face lowercase letters like $\mathbf x$ and $\mathbf y$ to denote sets of variables, $f,g$ for polynomials in $\mathbb F[\mathbf x,\mathbf y]$, uppercase typewriter alphabets $\mathtt C,\mathtt D$ for arithmetic circuits over $\mathbb F$, uppercase Roman alphabets $T,Q$ for the polynomials computed by the gates of a depth four circuit and $M,M_1,M_2$ for subsets of natural numbers. For $M\subseteq[3m]$, $\mathbf x_M:=\{x_i:i\in M\}$, $\mathbf y_M:=\{y_i:i\in M\}$. For $\mathbf z\subseteq\mathbf x_M\cup\mathbf y_M$ and $r\in\mathbb N$, $\mathbf z^{\le\infty}$ and $\mathbf z^{\le r}$ denote the set of all monomials in $\mathbf z$ variables and the set of all monomials in $\mathbf z$ variables with degree at most $r$ respectively. For $S\subseteq\mathbb F[\mathbf x,\mathbf y]$, $\dim\langle S\rangle$ denotes the dimension of the $\mathbb F$-linear span of $S$.
\paragraph{Support and degree of a monomial}
The support of a monomial $\eta$, denoted $\operatorname{Supp}(\eta)$, is the set of variables appearing in it. Also, for any $\mathbf z\subseteq\mathbf x\cup\mathbf y$ we will use $\deg_{\mathbf z}(\eta)$ to denote its degree in $\mathbf z$ variables. We will say that $\eta$ is $\mathbf z$-multilinear if the degree of every $\mathbf z$ variable in $\eta$ is at most one.
\subsection*{2.1 The complexity measure}
Throughout this section, we will assume that $m\in\mathbb N$ is as stated in the paragraph following Theorem1, $M\subseteq[3m]$, $|M|=m$, $f\in\mathbb F[\mathbf x_M,\mathbf y_M]$ and $S\subseteq\mathbb F[\mathbf x_M,\mathbf y_M]$. Note that the set $M$ is not fixed and will depend on the circuit under analysis. Before defining the measure, let us define the operations that make up the measure.
\begin{enumerate}
\item \textbf{Partial derivatives.} Let $\eta=x_1\cdots x_k$ be a monomial in $\mathbf x$ variables. Then, we define the partial derivative of $f$ with respect to $\eta$ as
\[
\frac{\partial f}{\partial\eta}:=\frac{\partial}{\partial x_1}\left(\frac{\partial}{\partial x_2}\left(\cdots\left(\frac{\partial f}{\partial x_k}\right)\right)\right).
\]
If the degree of $\eta$ is $k$, then $\partial f/\partial\eta$ is said to be a $k$-th order partial derivative of $f$. We denote by $\partial_{\mathbf x}^kf$ the set of all $k$-th order partial derivatives of $f$ taken with respect to multilinear monomials in $\mathbf x$ variables.
\item \textbf{The shift operation.} Let $\eta$ be a degree $\ell$ multilinear monomial in $\mathbf x_M$ variables. We say that the polynomial $\eta\cdot f$ is obtained by shifting $f$ by $\eta$. We denote by $\mathbf x_M^\ell f$ the set of polynomials obtained by shifting $f$ by all degree $\ell$ multilinear monomials in $\mathbf x_M$ variables and $\mathbf x_M^\ell S:=\{\mathbf x_M^\ell f:f\in S\}$.
\item \textbf{Multilinear projection.} We define a map $\pi_{\mathbf x}:\mathbb F[\mathbf x_M,\mathbf y_M]\to\mathbb F[\mathbf x_M,\mathbf y_M]$ with $\pi_{\mathbf x}(f)$ being the polynomial made up of exactly the $\mathbf x$-multilinear monomials of $f$. Formally, for a monomial $\eta$, $\pi_{\mathbf x}(\eta)=\eta$ if $\eta$ is $\mathbf x$-multilinear and 0 otherwise. The map is then linearly extended for arbitrary polynomials and $\pi_{\mathbf x}(S):=\{\pi_{\mathbf x}(f):f\in S\}$.
\item \textbf{A degree based projection.} For $i\in\mathbb N$ and $f\in\mathbb F[\mathbf x_M,\mathbf y_M]$, we define $[f]_i$ to be the polynomial made up of only those monomials of $f$ whose $\mathbf y$-degree is exactly $i$. Formally, for a monomial $\eta$, $[\eta]_i=\eta$ if $\deg_{\mathbf y}(\eta)=i$ and 0 otherwise. It is then linearly extended for arbitrary polynomials and $[S]_m:=\{[f]_m:f\in S\}$.
\item \textbf{An evaluation map.} For $\alpha\in\mathbb F$ and $\mathbf z\subseteq\mathbf x_M\cup\mathbf y_M$, we define a map $\sigma_{\mathbf z=\alpha}:\mathbb F[\mathbf x_M,\mathbf y_M]\to\mathbb F[\mathbf x_M\setminus\mathbf z,\mathbf y_M\setminus\mathbf z]$ with $\sigma_{\mathbf z=\alpha}(f)$ being obtained from $f$ by setting every variable in $\mathbf z$ to $\alpha$ and $\sigma_{\mathbf z=\alpha}(S):=\{\sigma_{\mathbf z=\alpha}(f):f\in S\}$.
\end{enumerate}
The operations given in1,2 and3 constitute the projected shifted partials measure [49]. In this work, we define and use the measure $\PSP_{M,k,\ell}$, which is obtained by augmenting the projected shifted partials measure with the operations in4 and5 as follows.
\begin{authorclaim}{Definition2 (The measure)}
For $m,k,\ell\in\mathbb N$, $M\subseteq[3m]$, $|M|=m$ and $f\in\mathbb F[\mathbf x_M,\mathbf y_M]$,
\[
\PSP_{M,k,\ell}(f):=\dim\left\langle\sigma_{\mathbf y_M=1}\left([\pi_{\mathbf x}(\mathbf x_M^\ell\partial_{\mathbf x}^kf)]_m\right)\right\rangle.
\]
\end{authorclaim}
\begin{authorclaim}{Observation3 (Sub-additivity of the measure)}
For any two polynomials $f,g\in\mathbb F[\mathbf x_M,\mathbf y_M]$, $\PSP_{M,k,\ell}(f+g)\le\PSP_{M,k,\ell}(f)+\PSP_{M,k,\ell}(g)$.
\end{authorclaim}
The above observation is easy to prove and we omit its proof here.
\subsection*{2.2 Some numerical estimates}
\begin{authorclaim}{Proposition4 (Estimating Binomial Coefficients)}
For any $n,k\in\mathbb N$, $k\le n$, $(n/k)^k\le\binom nk<(en/k)^k$.
\end{authorclaim}
\begin{authorclaim}{Proposition5 ([33,49])}
Let $a(n),f(n),g(n):\mathbb Z_{>0}\to\mathbb Z$ be integer values functions such that $(|f|+|g|)=o(a)$. Then, $\ln((a+f)!/(a-g)!)=(f+g)\ln(a)\pm O((f^2+g^2)/a)$.
\end{authorclaim}
